%!TEX root = ../template.tex
\typeout{NT FILE chapter-introduction.tex}%

\chapter{Introduction}
\label{cha:introduction}

% ===================================================================
\section{Context \& Motivation}
\label{sec:intro-context}
% ===================================================================

The Grande Panorama de Lisboa is a 23-metre azulejo tile panel, attributed
on stylistic grounds to the Spanish ceramic painter Gabriel del Barco and
generally dated to around 1700 \cite{gacpanorama,baroqueartpanorama}. It
depicts roughly 14\,km of the Tagus waterfront from Alg\'{e}s to Xabregas,
rendering palaces, churches, convents, and ordinary dwellings across more
than 150 identifiable structures, and survives today as the single most
complete pictorial record of Lisbon \emph{before} the 1755 earthquake,
which destroyed an estimated 85\% of what it depicts. A visitor standing
before this panel at the Museu Nacional do Azulejo (MNAz) today,
however, encounters very little beyond a single mounted label: a
physically rich artefact containing 150 or more buildings is, in
practice, uninterpretable without expert architectural knowledge, and
most visitors move past it in under two minutes. Mobile AR offers a way
to close that gap largely without cost to the institution: consumer
smartphones equipped with ARCore and ARKit can already overlay
interactive content onto a physical artefact, and recent bibliometric
work suggests this is no longer a marginal research direction. Cultural-heritage AR studies have been increasing year on
year through 2023 \cite{wang2025grand}, so the timing, if nothing else, seems right.

This thesis treats the Panorama as its central motivating case, but the
underlying problem is not specific to it. Any sufficiently large,
content-dense heritage wall raises the same questions: how to track a
visitor's position reliably across tens of metres, how to keep dozens or
hundreds of points of interest legible instead of overwhelming, and how
to serve visitors with very different backgrounds and available time.
To keep the generalisability claim honest, not aspirational,
a second surface is tested throughout: Chafariz Velho, a large,
tile-covered Pombaline-era fountain in Pa\c{c}o de Arcos. If anything
Chafariz Velho is the harder case. It is outdoor, circular, and subject
to shifting natural light, where the Panorama is indoor and
geometrically simpler, and (as \S\ref{sec:intro-institutional} explains)
it ended up as this thesis's primary development surface mostly by
accident of timing, not by original design. Both surfaces recur
throughout this report; neither is meant to be read as the ceiling of
what the resulting framework can support.

% ===================================================================
\section{Institutional Context \& Research Partnership}
\label{sec:intro-institutional}
% ===================================================================

MNAz is a high-traffic, high-value setting for this work: it was
Lisbon's most-visited national museum in 2024, with 297{,}203 recorded
entries \cite{observador2025mnaz}. Since 1 November 2025, however, the
museum has been closed to the public for renovation funded under
Portugal's Plano de Recupera\c{c}\~{a}o e Resili\^{e}ncia, with no
confirmed reopening date at the time of writing \cite{publico2025prr}.
No formal partnership agreement with MNAz was in place when this report
was prepared, and none is assumed by the plan that follows. Given this,
Chafariz Velho became the primary surface for the tracking tests,
mapping work, and early development reported in
Chapter~\ref{ch:proposed-work}. It is permanently and freely accessible
and needs no institutional sign-off at all. That is why it is
treated as the guaranteed setting for this thesis's evaluation
(\S\ref{sec:pw-evaluation}). Deployment and evaluation at MNAz itself
remain a goal pursued opportunistically, should the museum reopen with
enough time left before submission, but no research question,
contribution, or evaluation claim in this thesis depends on that
happening.

% ===================================================================
\section{Problem Statement}
\label{sec:intro-problem}
% ===================================================================

Four related gaps motivate this thesis, expanded with full justification
in \S\ref{sec:pw-problem-rq}. Visitors cannot currently identify
individual buildings or orient themselves within a panel as dense as the
Panorama; the panel fixes one historical moment with no mechanism
connecting it to the 1755 earthquake, the Pombaline reconstruction, or
the present-day city; passive observation of a static, label-only display
yields low dwell time and little retained understanding; and a single
fixed text cannot reasonably serve a child, a tourist, and an academic
historian at once. Together, these gaps describe a more general problem:
existing interpretation tools for large, content-dense heritage walls do
not scale to the density of information such walls actually carry.

% ===================================================================
\section{Research Questions}
\label{sec:intro-rqs}
% ===================================================================

This thesis is organised around one central question, decomposed into
five sub-questions in \S\ref{sec:pw-problem-rq}: \emph{how can a mobile
AR framework, built around an existing visual positioning service,
provide stable multi-POI spatial registration, historically layered
content, and visitor-adaptive personalisation across large-scale heritage
wall surfaces --- and what engagement, usability, and learning outcomes
does such a framework produce for real visitors?} The sub-questions
address, in turn, tracking configuration, interface scalability,
engagement patterns under free use, personalisation effectiveness, and
immediate knowledge gain, each scoped specifically to what this thesis's
evaluation plan can realistically measure.

% ===================================================================
\section{Main Contributions}
\label{sec:intro-contributions}
% ===================================================================

This thesis's primary contribution is a configured and empirically
tested mobile AR \emph{framework} for large-scale heritage walls. 
It is a config-driven architecture: a new wall is added as data, and it runs on an existing visual positioning service rather than a purpose-built tracking method.
Generalisability is supported directly by
running the same framework, unmodified, across three structurally
distinct surfaces, and demonstrated concretely by packaging the proven
core as an installable Unity template and measuring how quickly a new
wall can be integrated against it. Three further contributions follow
from this core. The first is a temporal visualisation mechanism,
generalised beyond any single wall's content: on the Panorama it layers
the depicted period against later historical states and the present
day, while the same underlying state-axis mechanism turns out to extend
to a non-historical dimension, such as time-of-day, on a structurally
different wall. The second is an adaptive visitor-profiling engine that
adjusts content depth, vocabulary, and suggested route length to a
small set of self-reported visitor profiles established during a short
onboarding step. The third is an in-situ evaluation, conducted
primarily at Chafariz Velho with realistic, street-recruited samples in
the order of twenty to thirty participants and extended opportunistically
to MNAz should it reopen in time, measuring usability, engagement
patterns, and immediate knowledge gain. This last piece is reported as
exploratory evidence; with fifteen to thirty participants, it cannot
support strong statistical generalisation.

% ===================================================================
\section{Document Structure}
\label{sec:intro-structure}
% ===================================================================

Chapter~\ref{ch:related-work} reviews digital heritage interpretation,
AR concepts and tracking techniques, AR applications in museums and at
large-scale surfaces specifically, and UX, personalisation, and
evaluation practice in heritage AR, closing with the explicit gap this
thesis addresses. Chapter~\ref{ch:proposed-work} develops the research
questions in full, introduces the Panorama, Chafariz Velho, and a third,
structurally distinct wall as case studies motivating a framework rather
than a single-installation design, sets out the Design Science Research
approach followed, presents the TileStories framework concept and its
phased feature set, justifies the technology stack and system
architecture, and details the evaluation plan. Chapter~\ref{ch:work-plan} reports work already completed, the
phased development timeline with its Gantt chart, and a risk analysis
covering the dependencies (chiefly museum access) that sit outside this
thesis's direct control.